\documentclass{article}
\usepackage{amsmath,amssymb}
\usepackage[margin=1in]{geometry}
\allowdisplaybreaks
\usepackage[UTF8]{ctex}
\begin{document}

\section*{全视野随机整数多商品资源分配策略模型}

\begin{align*}
\max_{\pi}\quad &\mathbb{E}\left\{\left.\sum_{t\in\mathcal{T}}c_tX_t^{\pi}(R_t,D_t)\,\right|\,R_1\right\}\\
\text{subject to}\quad
&x_t=X_t^{\pi}(R_t,D_t),\qquad t\in\mathcal{T},\\
&(x_t,R_{t+1})\in\mathcal{Y}(R_t,D_t),\qquad t\in\mathcal{T}.
\end{align*}
这里的\(\pi\)是一个确定性马尔可夫策略，它由一系列决策函数\(\{X_t^{\pi}:t\in\mathcal{T}\}\)所定义。每个函数都将当前资源状态以及观察到的需求结果映射到一个可行的决策。该策略的可行性适用于使用这些函数的各种状态和需求结果，而不仅仅适用于某个特定的需求实现情况。初始资源状态\(R_1\)已经给定；后续的资源状态则通过下面的决策和转换来生成。

\paragraph{选定的扩展和来源范围。}
目标和政策定义位于物理PDF文件的第6页。以下可行的替换方案将客户索引替换应用于第6-7页，同时将多周期旅行替换应用于第7页，从而对第5-6页的原始定义进行替换。该论文分别列出了这两种替换方案；它们在第17页的第6节中被明确说明可以同时使用。因此，旅行时间转换中的“加载移动”项被替换为针对所有客户的累计值。单周期转换和资源平衡则被替换，而不是被额外添加。来自(5)的整数域仍被保留，用于客户索引加载移动的情况。

\section*{可行的决策与资源转移}
对于每一个\(t\in\mathcal{T}\)，所选择的可行决策集是
\[
\mathcal{X}(R_t,D_t)=
\left\{x_t:\begin{aligned}
&\sum_{j\in\mathcal{I}}\left(x_{ijt}^{ke}+\sum_{d\in\mathcal{D}}x_{ijt}^{kdl}\right)=R_{itt}^{k}
&&i\in\mathcal{I},\ k\in\mathcal{K},\\
&\sum_{k\in\mathcal{K}}x_{ijt}^{kdl}\leq D_{ijt}^{d}
&&i,j\in\mathcal{I},\ d\in\mathcal{D},\\
&x_{ijt}^{ke}\in\mathbb{Z}_{+}
&&i,j\in\mathcal{I},\ k\in\mathcal{K},\\
&x_{ijt}^{kdl}\in\mathbb{Z}_{+}
&&i,j\in\mathcal{I},\ k\in\mathcal{K},\ d\in\mathcal{D}
\end{aligned}\right\}.
\]
相应的可行转换集合为：
\[
\mathcal{Y}(R_t,D_t)=
\left\{(x_t,R_{t+1}):\begin{aligned}
&R_{jt',t+1}^{k}=\sum_{i\in\mathcal{I}}\left(
\mathbf{1}_{\tau_{ij}}(t'-t)x_{ijt}^{ke}
+\mathbf{1}_{\tau_{ij}}(t'-t)\sum_{d\in\mathcal{D}}x_{ijt}^{kdl}
\right)+R_{jt't}^{k},\\[-1mm]
&\hspace{35mm}j\in\mathcal{I},\ k\in\mathcal{K},\ t'>t,\\
&x_t\in\mathcal{X}(R_t,D_t)
\end{aligned}\right\}.
\]
因此，所有可用的车辆都被分配到了空载或负载运输任务上；所有类型的车辆的负载运输任务都共享着每个起点、终点以及客户的需求限制。在转换过程中使用的指标，具体定义见第7页。
\[
\mathbf{1}_{y}(x)=
\begin{cases}
1,&x=y,\\
0,&\text{otherwise}.
\end{cases}
\]

\section*{索引、数据、决策以及状态}
\[
\begin{aligned}
\mathcal{T}&=\{1,\ldots,T\}&&\text{set of periods over which demands occur},\\
\mathcal{I}&=\text{set of locations},\\
\mathcal{K}&=\text{set of vehicle types},\\
\mathcal{D}&=\text{set of customers}.
\end{aligned}
\]
客户扩展将\(\mathcal{D}\)定义为客户集合；讨论和实验中也提到了客户类型。对于\(i,j\in\mathcal{I}\)、\(k\in\mathcal{K}\)、\(d\in\mathcal{D}\)和\(t\in\mathcal{T}\)：
\begin{itemize}
\item \(D_{ijt}^{d}\)是指客户\(d\)在时期\(t\)内提出的必须从\(i\)运输到\(j\)的随机数量的需求。
\item \(x_{ijt}^{ke}\)是关于在时期\(t\)内，从\(i\)到\(j\)区域，有数量未知的类型为\(k\)的车辆空驶的决策。
\item \(x_{ijt}^{kdl}\)是关于在时段\(t\)，为客户\(d\)从\(i\)运输到\(j\)的、类型为\(k\)的车辆数量的决策。
\item \(c_{ijt}^{ke}\)是相应空运的固定成本。
\item \(c_{ijt}^{kdl}\)是对应服务客户\(d\)的固定净利润。
\end{itemize}
在此公式中，运动成本与利润仍属于抽象数据；实验成本生成规则并未被明确设定。

在客户扩展模式下，第5页中的向量表示方式变为
\[
\begin{aligned}
D_t&=[D_{ijt}^{d}]_{i,j\in\mathcal{I},\ d\in\mathcal{D}},\\
x_t^e&=[x_{ijt}^{ke}]_{i,j\in\mathcal{I},\ k\in\mathcal{K}},&
x_t^l&=[x_{ijt}^{kdl}]_{i,j\in\mathcal{I},\ k\in\mathcal{K},\ d\in\mathcal{D}},\\
c_t^e&=[c_{ijt}^{ke}]_{i,j\in\mathcal{I},\ k\in\mathcal{K}},&
c_t^l&=[c_{ijt}^{kdl}]_{i,j\in\mathcal{I},\ k\in\mathcal{K},\ d\in\mathcal{D}},\\
x_t&=[x_t^e\mid x_t^l],&c_t&=[-c_t^e\mid c_t^l].
\end{aligned}
\]
特别是，政策目标中的奖励是
\[
c_tx_t=\sum_{i,j\in\mathcal{I}}\sum_{k\in\mathcal{K}}
\left(-c_{ijt}^{ke}x_{ijt}^{ke}+\sum_{d\in\mathcal{D}}c_{ijt}^{kdl}x_{ijt}^{kdl}\right).
\]

\paragraph{需求假设和服务解释。}
第4页假设了时间上的独立性，原始表述为：\(D_{ijt}\)和\(D_{i'j't'}\)对于\(t\ne t'\)来说是独立的。第6页利用这一假设来证明确定性马尔可夫策略的可行性。客户扩展部分引入了\(D_{ijt}^{d}\)，但没有单独重新表述客户之间的独立性公式；也就是说，同一时期内不同客户之间不存在独立性。在核心模型中，需求分配是抽象的。某个地点的需求只能由该地点的车辆来满足。所有客户的需求都必须在应用程序中得到满足，但那些无法由内部车辆满足的需求则被视为丢失，由租赁车辆来提供。负载移动净利润是指内部服务与租赁服务之间的增量利润。因此，打印出来的内部需求条件实际上是一个不等式，而不是等式。

\paragraph{输入资源的状态和旅行时间。}
在第7页中，\(R_{it't}^{k}\)表示在时间\(t\)时，类型为\(k\)的车辆前往位置\(i\)的数量，这些车辆将在时间\(t'\)到达那里；定义中提到了\(i\in\mathcal{I}\)、\(k\in\mathcal{K}\)、\(t\in\mathcal{T}\)和\(t'\in\mathcal{T}\)。组件\(R_{itt}^{k}\)则负责提供在时间\(t\)时可供使用的车辆。固定数据\(\tau_{ij}\)表示从\(i\)到\(j\)的行驶时间。第5.3节，第14页，明确给出了多周期状态向量。
\[
R_t=[R_{it't}^{k}]_{i\in\mathcal{I},\ k\in\mathcal{K},\ t\leq t'<t+\tau_{\max}},
\]
其中，\(\tau_{\max}\)表示可能的最大旅行时间。

\paragraph{印刷版索引资料。}
第7页的定义中提到了\(t'\in\mathcal{T}\)，而其过渡范围则被标记为\(t'>t\)；第14页则给出了状态向量的范围\(t\leq t'<t+\tau_{\max}\)，并且没有在\(T\)处进行截断。这些被标记的範圍都保留了下来。顯示的轉換過程中也保留了\(R_{jt't}^{k}\)這個項目。對於顯示狀態向量範圍之外的元件，並沒有采用零填充的規則，同樣也沒有對移動行為施加額外的限制。這些規則都是為了讓這些範圍能夠保持一致而設定的。

\section*{链接的精确值函数方程}
第 6 页上的精确递归过程使用了相同的可行转移集合，并且采用了上述选定的替换方案。
\begin{equation}
V_t(R_t)=\mathbb{E}\left\{\left.
\max_{(x_t,R_{t+1})\in\mathcal{Y}(R_t,D_t)}
\bigl\{c_tx_t+V_{t+1}(R_{t+1})\bigr\}
\,\right|\,R_t\right\}.
\tag{6}
\end{equation}
第8页中，位于(7)之前的限定关系为：
\[
V_t(R_t)=\mathbb{E}\{V_t(R_t,D_t)\mid R_t\},
\]
其中
\begin{equation}
V_t(R_t,D_t)=
\max_{(x_t,R_{t+1})\in\mathcal{Y}(R_t,D_t)}
\bigl\{c_tx_t+V_{t+1}(R_{t+1})\bigr\}.
\tag{7}
\end{equation}
该页面图像使用相同的普通符号\(V_t\)来表示单参数期望值函数和双参数需求条件函数；它们的参数区分了各自的功能。打印的递归过程并未为\(V_{T+1}\)提供终止值方程；也没有添加任何相关内容。蒙特卡洛近似（8）、价值函数近似、后验确定性程序以及更新算法均不属于本政策制定的范畴。

\end{document}
